% =========================================================
% Computational Linear Algebra --- Introduction to Matrices
% (Strang \S{}1.3 / Johnston \S{}1.3)
% =========================================================

\section{Introduction to Matrices}

% ---------------------------------------------------------
\subsection*{Matrices and Matrix--Vector Multiplication}

\begin{propositionbox}{Definition (Matrix)}
An $m \times n$ \emph{matrix} $A$ is a rectangular array of real numbers with
$m$ rows and $n$ columns:
\[
  A = \begin{pmatrix}
    A_{11} & A_{12} & \cdots & A_{1n} \\
    A_{21} & A_{22} & \cdots & A_{2n} \\
    \vdots & \vdots & \ddots & \vdots \\
    A_{m1} & A_{m2} & \cdots & A_{mn}
  \end{pmatrix}.
\]
The entry in row $i$, column $j$ is denoted $A_{ij}$ or $[A]_{ij}$.
We write $A \in \mathbb{R}^{m \times n}$.
\end{propositionbox}

\begin{propositionbox}{Definition (Matrix--Vector and Matrix--Matrix Multiplication)}
If $A \in \mathbb{R}^{m \times n}$ and $\mathbf{x} \in \mathbb{R}^n$, then
$A\mathbf{x} \in \mathbb{R}^m$ is defined by
\[
  (A\mathbf{x})_i = \sum_{j=1}^n A_{ij} x_j, \qquad i = 1, \ldots, m.
\]
Equivalently, $A\mathbf{x}$ is the linear combination of the columns of $A$
with coefficients from $\mathbf{x}$:
\[
  A\mathbf{x} = x_1 \mathbf{a}_1 + x_2 \mathbf{a}_2 + \cdots + x_n \mathbf{a}_n.
\]
If $A \in \mathbb{R}^{m \times n}$ and $B \in \mathbb{R}^{n \times p}$, their product
$AB \in \mathbb{R}^{m \times p}$ has entries
\[
  (AB)_{ij} = \sum_{k=1}^n A_{ik} B_{kj}.
\]
\end{propositionbox}

\begin{remark*}[Matrix multiplication is not commutative]
In general $AB \neq BA$, even when both products are defined.
Order matters: applying transformation $B$ first, then $A$, is written $AB$.
\end{remark*}

% ---------------------------------------------------------
\subsection*{The Transpose}

\begin{propositionbox}{Definition (Transpose)}
The \emph{transpose} of $A \in \mathbb{R}^{m \times n}$ is the matrix
$A^T \in \mathbb{R}^{n \times m}$ defined by $(A^T)_{ij} = A_{ji}$.
Rows of $A$ become columns of $A^T$.
Key properties:
\begin{itemize}
  \item $(A^T)^T = A$
  \item $(AB)^T = B^T A^T$
  \item $\mathbf{v} \cdot \mathbf{w} = \mathbf{v}^T \mathbf{w}$ (dot product as matrix multiplication)
\end{itemize}
\end{propositionbox}

% ---------------------------------------------------------
\subsection*{Matrix Form of a Curve}

\begin{remark*}[NURBS connection: matrix form of curves and surfaces]
The power basis curve $C(u) = \sum_{i=0}^n \mathbf{a}_i u^i$ can be written as
\[
  C(u) = [\mathbf{a}_0 \;\; \mathbf{a}_1 \;\; \cdots \;\; \mathbf{a}_n]
  \begin{pmatrix} 1 \\ u \\ \vdots \\ u^n \end{pmatrix}
  = A \,\mathbf{u}(u),
\]
where $A$ is the matrix of coefficient vectors and $\mathbf{u}(u)$ is the monomial
vector.
A B\'{e}zier curve has the equivalent matrix form $C(u) = [u^i]^T M [P_i]$,
where $M$ is the $(n+1)\times(n+1)$ change-of-basis matrix converting
Bernstein basis to power basis.
A tensor product surface takes the form
$S(u,v) = [f_i(u)]^T [b_{i,j}] [g_j(v)]$,
a matrix sandwiched between two basis-function vectors.
All of these representations rely on matrix--vector multiplication as
defined above.
\end{remark*}
